\documentclass[11pt]{article}

\usepackage[a4paper,margin=1in]{geometry}
\usepackage{booktabs}
\usepackage{enumitem}
\usepackage{hyperref}
\usepackage{xcolor}

\hypersetup{
  colorlinks=true,
  linkcolor=blue,
  urlcolor=blue
}

\setlist[itemize]{noitemsep, topsep=3pt}
\setlist[enumerate]{noitemsep, topsep=3pt}

\title{Futoshiki VifeAgent system design}
\author{}
\date{}

\begin{document}
\maketitle

\section{overview}

This chapter describes the current structure of the Futoshiki VifeAgent
repository. The code is built around one puzzle type, but it is used in several
ways: solving puzzle files, running experiments, generating benchmark cases,
validating results, and showing solver behaviour in an interactive UI.

The repository compares brute force search, forward chaining, backward
chaining, A* search, constraint propagation, and Z3-based validation. These
parts share the same puzzle model. That is why the code is split into domain,
logic, solver, UI, benchmark, and CLI packages instead of being organised around
only one workflow.

There are three user-facing workflows:

\begin{itemize}
  \item a pygame UI for playing puzzles, inspecting knowledge bases, and
        watching solvers run step by step;
  \item a solver CLI for solving one puzzle file or a folder of puzzle files;
  \item benchmark tools for generating puzzles, validating the corpus, running
        solvers, and plotting results.
\end{itemize}

The basic design choice is to keep the puzzle representation separate from the
tools that use it. A \texttt{Puzzle} should not know whether pygame is showing
it, a solver is working on it, Z3 is checking it, or the formatter is writing
it to a file. The same domain object moves through all of those paths.

The repository is organised by responsibility:

\begin{verbatim}
src/
  core/                 Puzzle, Parser, Formatter
  constraints/          inequality constraints and AC-3 support
  fol/                  predicates, axioms, CNF and Horn knowledge bases
  inference/            forward and backward chaining engines
  heuristics/           A* heuristic functions
  search/               A* state and search engine
  solver/               BaseSolver and concrete solver classes
  models/               board state, game state, puzzle repository
  app/                  pygame application controller and input handling
  ui/                   renderers, layout, theme
  benchmark/            benchmark runner, generator, validator, plots
  futoshiki_vifeagent/  package entry points and CLI wrappers
\end{verbatim}

A workflow-based layout was possible. We could have used separate folders for
\texttt{play}, \texttt{solve}, and \texttt{benchmark}. That would read nicely
at first, but parsing, formatting, solver selection, and validation would start
appearing in several places. The repository uses more packages instead, with a
clearer job for each one.

\section{architecture and dependency direction}

The dependency direction is simple. Domain and logic packages sit near the
bottom. Application code sits near the top.

\begin{verbatim}
UI / CLI / benchmark
        |
      solver
        |
search, heuristics, inference
        |
   fol, constraints
        |
       core
\end{verbatim}

The lower packages do not import the UI, benchmark runner, or CLI. For example,
\texttt{core} can be tested without pygame; \texttt{inference} can be tested
without the puzzle generator; \texttt{solver} can be used by the CLI or by the
UI without changing its public contract.

The result is a useful boundary between "what the puzzle is" and "what we do
with the puzzle." Core puzzle data and reasoning code stay independent from the
runtime surfaces. The UI, CLI, and benchmark tools adapt that behaviour for
different users, but they do not define the puzzle model itself.

\section{core domain}

\subsection{Puzzle as the problem definition}

\texttt{Puzzle} is the central domain object. It stores:

\begin{itemize}
  \item the grid size \texttt{N};
  \item a NumPy grid where \texttt{0} means empty and \texttt{1..N} are values;
  \item horizontal and vertical \texttt{InequalityConstraint} objects;
  \item cached lists of given cells and empty cells;
  \item lookup maps for constraints between neighbouring cells.
\end{itemize}

\texttt{Puzzle} represents the problem, not the whole application state.
Solvers receive a puzzle and return either a solved puzzle or \texttt{None}.
The UI does not edit the original puzzle directly. It wraps the puzzle in
\texttt{Board}, which keeps player edits, notes, selected cell state, errors,
and undo history.

\begin{table}[h]
\centering
\begin{tabular}{lll}
\toprule
Class & File & Role \\
\midrule
\texttt{Puzzle} & \texttt{core/puzzle.py} & Immutable puzzle definition \\
\texttt{Parser} & \texttt{core/parser.py} & Reads puzzle text files \\
\texttt{Formatter} & \texttt{core/formatter.py} & Renders puzzles for output \\
\texttt{Board} & \texttt{models/board.py} & Mutable UI board state \\
\texttt{SearchState} & \texttt{search/state.py} & A* node state \\
\bottomrule
\end{tabular}
\end{table}

We could have made \texttt{Puzzle} mutable and let every subsystem edit it.
That would reduce the number of classes, but it would blur the meaning of the
grid. A grid could be the original puzzle, a player's partial attempt, or a
temporary solver state. Keeping those meanings separate makes the code easier
to reason about.

\section{class hierarchy}

Inheritance is used only where a caller really needs interchangeable
implementations. Most objects are plain data objects or service classes.

\subsection{solver hierarchy}

\begin{verbatim}
BaseSolver
  BruteForceSolver
  ForwardChaining
  BackwardChaining
  BacktrackingForwardChaining
  ForwardThenBackwardChaining
  AStarSolver
\end{verbatim}

All solvers implement the same contract:

\begin{verbatim}
solve(puzzle, on_step=None) -> (Puzzle | None, Stats)
get_name() -> str
\end{verbatim}

The \texttt{on\_step} callback is optional. The CLI and benchmark runner ignore
it. The UI uses it to animate assignments and backtracking.

\subsection{heuristic hierarchy}

\begin{verbatim}
BaseHeuristic
  EmptyCellHeuristic
  DomainSizeHeuristic
  MinConflictsHeuristic
  AC3Heuristic
\end{verbatim}

\texttt{AStarSolver} receives a heuristic object instead of hard-coding one
estimate. The A* engine stays stable, while experiments can swap in different
heuristics.

\subsection{constraint and renderer hierarchy}

\begin{verbatim}
BaseConstraint
  RowUniqueness
  ColUniqueness
  InequalityConstraint

BaseRenderer
  CompositeRenderer
  GridRenderer
  ConstraintRenderer
  HudRenderer
\end{verbatim}

\texttt{InequalityConstraint} is part of the puzzle definition. Row and column
constraints mainly support AC-3 domain pruning. The renderer hierarchy is
similarly shallow: the composite renderer controls draw order, and the concrete
renderers draw the grid, inequality symbols, and HUD.

A deeper object-oriented framework would make the hierarchy look more formal,
but it would add little value here. The hierarchy is small enough to read
directly.

\section{FOL representation}

The FOL package is the logical model of the repository. It translates a
\texttt{Puzzle} into predicates, clauses, and knowledge bases that other code
can use. It does not run solvers, draw the UI, read benchmark folders, or choose
an algorithm.

Its interface is small: a generator receives a \texttt{Puzzle} and returns a
knowledge-base object. The caller decides whether to use that object for
display, counting, forward chaining, or backward chaining. Runtime concerns
stay outside the FOL layer.

\subsection{Predicates and literals}

\texttt{predicates.py} defines the vocabulary used by the logic layer. Common
predicates include:

\begin{itemize}
  \item \texttt{Val(r,c,v)}: cell \texttt{(r,c)} has value \texttt{v};
  \item \texttt{NotVal(r,c,v)}: cell \texttt{(r,c)} cannot have value
        \texttt{v};
  \item \texttt{Given(r,c,v)}: the puzzle provides a fixed value;
  \item \texttt{Less}, \texttt{Geq}, and \texttt{Diff}: helper relations for
        inequality and uniqueness rules.
\end{itemize}

These predicates are the interface between the puzzle domain and the reasoning
code. They give the rest of the repository stable names for cell values, fixed
clues, inequality relations, and value exclusions. The tradeoff is that the
predicates are lower-level than the puzzle itself. A caller has to know what
\texttt{Val} or \texttt{NotVal} means, but the inference engines get a clean
input shape that is not tied to pygame or files.

\subsection{Axioms and CNF knowledge base}

\texttt{axioms.py} contains the rule schemas for the puzzle. The schemas cover
the standard Futoshiki responsibilities:

\begin{itemize}
  \item each cell has a value in the domain;
  \item a cell cannot contain two different values;
  \item rows and columns cannot repeat values;
  \item given cells must keep their value;
  \item inequality constraints must hold between neighbouring cells.
\end{itemize}

\texttt{CNFGenerator} grounds these schemas for a concrete puzzle and stores
the result in \texttt{CNFClauseKnowledgeBase}. The responsibility of this path
is representation, not solving. It gives the UI and report code a consistent
view of the logical constraints behind a puzzle.

The tradeoff is readability versus direct execution. CNF is a familiar and
compact way to display constraints, but the chaining engines do not use it
directly. Keeping CNF separate avoids a forced universal clause format.

\subsection{Horn knowledge bases}

The repository has two Horn clause generators:

\begin{itemize}
  \item \texttt{HornClauseGenerator}, used by backward chaining. It creates a
        goal-oriented knowledge base. Its public contract is
        \texttt{generate(puzzle, domains=None, use\_cell\_domains=False)}
        \(\rightarrow\) \texttt{HornClauseKnowledgeBase}.
  \item \texttt{HornClauseGenerator2}, used by forward chaining. It creates
        a propagation-oriented knowledge base. Its public contract is
        \texttt{generate(puzzle)} \(\rightarrow\)
        \texttt{HornClauseKnowledgeBase}.
\end{itemize}

Forward chaining needs local rules that can derive intermediate facts.
Backward chaining needs a goal-oriented encoding for complete assignments. A
single shared Horn encoding would either be too coarse for propagation or too
fragmented for goal-directed search. The repository therefore keeps two Horn
generators with the same domain meaning but different interfaces.

\section{inference design}

The inference package owns the rule-application engines. It does not know about
files, pygame state, benchmark output, or puzzle formatting. It receives Horn
clauses, facts, goals, and substitutions. It returns derived facts or
substitutions.

The boundary is straightforward. FOL generators decide how a Futoshiki puzzle
becomes a Horn knowledge base. Inference engines decide how Horn data is used.
Solvers sit above both layers and handle the path from puzzle to knowledge base
to solved puzzle.

\subsection{Forward chaining}

\texttt{ForwardChainingEngine} receives a list of \texttt{HornClause} rules and
a list of initial \texttt{Literal} facts. Its main method,
\texttt{run()}, returns the final set of known facts. The engine is responsible
for applying the rules to its fact store; it is not responsible for deciding
which facts count as a solved Futoshiki board.

\texttt{ForwardChainingEngine} consumes Horn rules and initial facts, then
produces a saturated fact set. The \texttt{ForwardChaining} solver interprets
those facts as puzzle assignments.

\subsection{Backward chaining}

\texttt{BackwardChainingEngine} receives a
\texttt{HornClauseKnowledgeBase}. Its main methods accept one or more goal
literals and return substitutions: \texttt{prove(goal)} handles one goal,
\texttt{prove\_all(goals)} returns the first satisfying substitution, and
\texttt{prove\_all\_solutions(goals)} yields satisfying substitutions.

\texttt{BackwardChainingEngine} runs goal-directed inference over Horn clauses
and returns satisfying substitutions through a generator-based interface. It is
separate from puzzle-specific rule generation, so it can be tested without the
Futoshiki encoding.

\subsection{Why inference is separate from solvers}

The solvers decide which knowledge base to build and how to translate the
engine result back into a \texttt{Puzzle}. The engines decide only how to
consume Horn data. This gives the code a concrete testing boundary: engine
tests can use small artificial Horn programs, while solver tests can check
puzzle-level behaviour.

We could have put the inference loop directly inside each solver. That would
reduce the number of files, but it would mix Futoshiki-specific rule
generation, result conversion, and general inference mechanics. The code has
more small classes this way, but each class has a clearer reason to change.

\section{solvers and search}

The solver layer is the common boundary used by the UI, CLI, and benchmark
runner. Every solver returns a solution and a \texttt{Stats} object. The
selected algorithm changes, but the caller does not.

\begin{table}[h]
\centering
\begin{tabular}{lll}
\toprule
Key & Class & Main idea \\
\midrule
\texttt{brute\_force} & \texttt{BruteForceSolver} & enumerate possible assignments \\
\texttt{forward\_chaining} & \texttt{ForwardChaining} & derive facts until stable \\
\texttt{backward\_chaining} & \texttt{BackwardChaining} & prove a solution goal \\
\texttt{backtracking\_forward\_chaining} & \texttt{BacktrackingForwardChaining} & propagation plus search \\
\texttt{forward\_then\_backward\_chaining} & \texttt{ForwardThenBackwardChaining} & FC first, BC fallback \\
\texttt{astar\_h1} & \texttt{AStarSolver} & A* with empty-cell heuristic \\
\texttt{astar\_h2} & \texttt{AStarSolver} & A* with domain-size heuristic \\
\texttt{astar\_h3} & \texttt{AStarSolver} & A* with conflict heuristic \\
\texttt{astar\_h4} & \texttt{AStarSolver} & A* with AC-3-pruned domains \\
\bottomrule
\end{tabular}
\end{table}

\subsection{Why several solvers exist}

The repository exposes several solver families behind one shared interface
because the report compares different reasoning and search strategies. The UI,
CLI, and benchmark runner can use the same integration path while swapping the
actual solver.

The cost is extra registry work. The CLI, UI, and benchmark path each need a
way to map user-facing solver keys to solver objects. The shared
\texttt{BaseSolver} interface keeps that manageable.

\subsection{Stats as a shared result shape}

All solvers return a shared \texttt{Stats} structure, so the UI, CLI, and
benchmark pipeline can report results in the same format. The fit is imperfect
because different algorithms naturally expose different metrics, but one shared
shape keeps the reporting layer simple.

\section{UI design}

The pygame UI uses three main objects:

\begin{itemize}
  \item \texttt{GameApplication}, which owns the main loop;
  \item \texttt{GameState}, which owns mutable UI state;
  \item \texttt{CompositeRenderer}, which draws the screen.
\end{itemize}

The UI has three modes:

\begin{table}[h]
\centering
\begin{tabular}{lll}
\toprule
Mode & Purpose & Main data \\
\midrule
\texttt{PLAY} & manual solving & board, notes, undo, errors \\
\texttt{SOLVE} & solver animation & step queue, solver cells, stats \\
\texttt{KB} & knowledge-base view & CNF clauses, scroll state, hover state \\
\bottomrule
\end{tabular}
\end{table}

\subsection{Single owner of UI state}

All mutable UI state lives in \texttt{GameState}. That object is large, but it
has one advantage: the main thread owns the UI state. Solver workers do not
write to it directly. They push \texttt{SolveStep} objects into a queue, and the
main loop consumes the queue during rendering.

We could have let solver threads update the board directly. That would be
shorter, but it would make pygame state harder to protect. The queue-based
model is slightly more code, but thread ownership is clearer.

\subsection{Renderer split}

Rendering is split into grid, constraint, and HUD renderers. Each renderer owns
one part of the screen. The composite renderer owns draw order, so the rest of
the application does not need to know which renderer runs first.

The UI is not fully model-view-controller. That would be more formal, but it
would also add more layers. The structure is simple enough for the project size
and still keeps input, state, and rendering mostly separate.

\section{generator and validator}

The generator and validator support the benchmark workflow. They are not part
of the solver interface because their job is different: produce and verify test
data.

\subsection{Generator}

\texttt{FutoshikiGenerator} creates benchmark puzzles in three phases:

\begin{enumerate}
  \item build a full valid grid with recursive backtracking;
  \item add inequality constraints using the completed solution;
  \item remove clues while preserving a unique solution.
\end{enumerate}

During generation, search heuristics and Z3-based uniqueness checks keep puzzle
construction practical and reliable.

\subsection{Validator}

\texttt{validator.py} parses benchmark files, checks obvious row, column, and
inequality conflicts, builds a Z3 model, and counts solutions. It can also
compare a puzzle against an expected solution file.

Z3 is used here because benchmark validation should not trust the same solver
being measured. If a solver contains a bug, the validator still needs to catch
wrong output. The cost is an extra dependency and a second modelling style, but
the independence is useful for benchmark checks.

\section{CLI design}

The package exposes a small command tree:

\begin{itemize}
  \item \texttt{futoshiki}, the top-level command;
  \item \texttt{futoshiki solver}, the direct solver tool;
  \item \texttt{futoshiki benchmark}, the benchmark tool group;
  \item \texttt{futoshiki ui}, the pygame launcher.
\end{itemize}

The command split is:

\begin{verbatim}
uv run futoshiki solver <file-or-folder> [--solver KEY] [--output-dir DIR]
uv run futoshiki benchmark run --solver KEY [--benchmark-root DIR] [--max-n N]
uv run futoshiki benchmark generate [--benchmark-root DIR]
uv run futoshiki benchmark validate <path> [--expected-dir DIR]
uv run futoshiki benchmark visualize [--data-dir DIR] [--output-dir DIR] [--latex-dir DIR]
uv run futoshiki ui
\end{verbatim}

\texttt{futoshiki solver} has one job. It receives either one \texttt{.txt}
puzzle file or a folder of \texttt{.txt} files. For one file, it prints the
solved board to standard output. For a folder, it writes
\texttt{\_solution.txt} files into an output directory. The command stays close
to the normal solver workflow: give it puzzle input, choose a solver if needed,
and get solved boards back.

\texttt{futoshiki benchmark} is a dispatcher for benchmark tools. It does not
copy the argument parsers from \texttt{benchmark.py}, \texttt{generator.py},
\texttt{validator.py}, or \texttt{visualize.py}. It forwards the remaining
arguments to the selected tool, and that tool handles its own options, help
text, and validation.

The tradeoff is a small routing layer. Adding a new benchmark tool means adding
the tool implementation and registering it in the benchmark command table. In
return, the solver command stays simple, and benchmark scripts keep the
interfaces they already expose.

\section{design summary}

\begin{table}[h]
\centering
\begin{tabular}{p{0.28\linewidth}p{0.32\linewidth}p{0.32\linewidth}}
\toprule
Choice & Benefit & Cost \\
\midrule
Small \texttt{Puzzle} object & shared by parser, solvers, UI, and tests & needs wrapper objects for mutable workflows \\
One \texttt{BaseSolver} interface & UI, CLI, and benchmark can swap solvers & one \texttt{Stats} shape cannot fit every algorithm perfectly \\
Separate FOL and inference layers & rules and engines can be tested separately & more setup code and more objects \\
Multiple solver families & supports comparison across AI/search strategies & registry and reporting code must stay in sync \\
Worker thread for solver animation & UI remains responsive & queue and stop-event plumbing \\
Z3 for generation and validation & independent uniqueness checks & extra dependency and separate modelling style \\
Direct solver CLI, benchmark tool router & solver use is short, benchmark tools keep their own interfaces & new benchmark tools must be registered in the router \\
\bottomrule
\end{tabular}
\end{table}

\end{document}
